\section{Evaluation}
\label{sec:eval}

We evaluate the system on three corpora and a separate routing audit, differing in origin, language
and scale. The evaluation follows the management levels of
Table~\ref{tab:levels} rather than the order in which the system was built:
first the collection level, where documents are parsed, classified and
catalogued (Section~\ref{sec:e1}), then the document level, where the
collection is kept consistent through lifecycle operations
(Section~\ref{sec:e2}), and finally the lower tiers of the access path, where
retrieval quality (Section~\ref{sec:e4}), answer-layer checks per question and
abstention (Section~\ref{sec:e5}, the paper's end-to-end evidence), a
deployment-layer replication through the public API (Section~\ref{sec:sysarms}),
verifier failure (Section~\ref{sec:e6}), non-functional behaviour
(Section~\ref{sec:e7}), the gate and search-scope ablations
(Section~\ref{sec:e8}), and the content-layer optimisation arms with the
ladder-versus-flat-and-precise contrast (Section~\ref{sec:ladder-opt}) are
measured. The collection-selection audit of
Section~\ref{sec:routing-audit} sits between the first two groups because it
exercises the tier that decides which collection is consulted. Numbers are
reported against the corpus that produced them, never merged; the campaign
inventory is supplementary. Table~\ref{tab:kb-checks} maps each management
function to its test and observed outcome; Table~\ref{tab:retrieval} reports
retrieval quality and Table~\ref{tab:e5-answers} the per-question answer
checks.

\subsection{Experimental setup}
\label{sec:eval-setup}

\textbf{Corpora.} (i)~\emph{Industrial demonstration corpus}: five libraries
hold 42 documents prepared as a demonstration instance, not customer data; one
library contains twenty near-duplicate
equipment-telemetry reports sharing prefixes, vocabulary and structure---a
deliberate stressor for near-duplicate detection and confuser
density~\cite{manku-near-dup-2007}; sixteen gold questions and four
out-of-corpus probes ran against it. (ii)~\emph{arXiv corpus}: 100 submitted
papers spanning about 25 domains; its campaign record contains 322 stored
parts across five libraries after the 30{,}000-character split. At the
2026-10-13 snapshot
79 papers were retrievable and 21 were not. Ten gold questions (BQ01--BQ10) target papers referenced by those
papers, keyed by arXiv identifier; the running instance then held 19 libraries. (iii)~\emph{Cross-lingual
ingestion corpus}: 400 Wikipedia pages per language (English, Chinese,
Japanese), 40 extractive cloze queries each with verbatim in-page answers;
an ingestion benchmark with no task-level retrieval-quality comparison.
(iv)~\emph{Routing
audit}: a separate ten-domain corpus of 138 documents and 879 stored parts,
with 72 paired in-domain questions and four out-of-domain probes per arm;
two cross-placeholder probes are excluded. Its route-selection outcome is
reported separately and is not merged with the other corpora's answer or
retrieval metrics.

\textbf{Questions and gold annotations.} The sixteen industrial questions were
authored, not sampled from logs; each identifies one or two gold documents by
a unique filename fragment and one to three fact markers. Q16 includes a token
also present in its gold filename, so a filename-ranking retriever could match
it without reading content; we retain this limitation and report a sensitivity
check in Section~\ref{sec:e4}. The arXiv answer grader
checks gold-ID presence in the answer or trace plus a word-boundary keyword
rule that counts a cell as covered when it matches at least $\lceil 3n/5\rceil$
of the $n$ annotated terms; it does not assess semantic support. Citation
markers are counted separately by a deterministic pattern and do not assess
source resolution or claim-to-source entailment.

\textbf{Metric families.} Retrieval uses Hit@5, Recall@5 and MRR@5, following
passage-ranking conventions~\cite{craswell-msmarco-2021}; Hit@5 and Recall@5
are document-level over the top five returned documents, and MRR@5 is $1/r$
for the first gold at rank $r\le 5$ and zero otherwise.
nDCG@5~\cite{jarvelin-cumulated-gain-2002} was computed only inside one
exploratory arm whose protocol was superseded before the reported campaigns;
all reported arms use Hit@5, Recall@5 and MRR@5. The routing audit reports gold-domain
knowledge-base choice accuracy with Wilson 95\% intervals and an exact paired
McNemar test. Its separate gold-string containment check counts an exact gold
string in the evidence pack and is not an answer-quality metric. The end-to-end
"gold-hit" measure checks
whether the answer or recorded retrieval trace contains the gold identifier;
keyword coverage is a deterministic lexical proxy. Citation counts record
whether the answer contains a part, section or table marker under the stated
pattern. These checks support traceability and marker-presence claims, not
semantic faithfulness or claim-to-source support, for which model-based
measures such as RAGAS~\cite{es-ragas-2024} or ARES~\cite{saad-falcon-ares-2024}
are designed. The third metric family covers wall-clock time, agent turns,
tool calls, cost and whether refusals were appropriate.

\textbf{Controls.} The deterministic retrieval lanes were re-run and compared
field by field, latency excluded; the 27-cell stability suite
(Section~\ref{sec:e7}) and the eighteen-arm honesty battery
(Section~\ref{sec:sysarms}) carry reruns of interrupted or failed runs rather
than independent repeated observations (the battery has one final observation
per arm), whereas the 30-cell answer grid of
Section~\ref{sec:e5} is a single run per arm on one model, recorded in the run
timeline as claude-opus-4-8[1m] with seed~0 and harness prompt v4-neutral, and
is reported as such, as is its controlled replication. The fixed retry
policy (one retry for
an empty zero-turn response or transient gateway timeout) fired twice,
succeeding both times. Before measurement we repaired two corpus-integrity
faults through the platform's own operations: one document missing from its
vector collection was re-indexed, and one library whose collection had been
lost was rebuilt, the same failure class that Section~\ref{sec:e2} measures. The tool-layer
retrieval lanes, grading rule, audits, duplicate adjudication and cross-lingual
ingestion pipeline are deterministic; agent lanes and text-derived results are
not. Tables~\ref{tab:kb-checks} and~\ref{tab:retrieval-design} summarise the
design and headline results of the management-layer and of the retrieval and
answer experiments respectively; the subsections that follow report each in
full.

\begin{table*}[t]
\centering
\footnotesize
\setlength{\tabcolsep}{4pt}
\caption{Functional checks on the management layer. Each row names the
function, the test and its scale, the observed outcome, and the subsection
that reports it. The numbers are those of the cited campaigns; none is
re-estimated. A passing check means the stated property was observed on the
stated scale, not that the function is exhaustively verified.}
\label{tab:kb-checks}
\begin{tabular}{@{}>{\raggedright\arraybackslash}p{0.115\textwidth}>{\raggedright\arraybackslash}p{0.265\textwidth}>{\raggedright\arraybackslash}p{0.39\textwidth}>{\raggedright\arraybackslash}p{0.10\textwidth}@{}}
\toprule
Function & Test and scale & Observed outcome & Where\\
\midrule
Parse and ingest & 100-paper campaign; state-file reconciliation & 91 names in the state file; 79 retrievable at the 2026-10-13 snapshot and 21 not (twelve recorded, nine unrecorded); no per-paper parse status & E1\\
Parse quality & 50-paper, 26-domain campaign & 50/50 parsed and verified, all probes ok; 202 tags; 41 tag flags negative without blocking ingestion & E1\\
Structural split & 50 papers at the 30{,}000-character tier & 165 stored parts (9 whole papers, 156 split parts) in 2{,}672\,s; storage 50/50; descriptions and tags 165/165 & E1\\
Near-duplicate review & 322-part corpus, 0.90 threshold & 35 candidate groups (0.90--0.975), zero true duplicates, nothing deleted & E2\\
Content placement & 100-paper campaign; 400 pages per language & 322 parts recorded for the paper campaign; cross-lingual membership 0.875 (en), 1.0 (zh, ja) & E1--E2\\
Content audit and three-source consistency & 322-part content audit; 33-part storage sample & 322/322 content; 33/33 sampled storage checks; separate disposable-library rename defect fixed by re-index & E2\\
Lifecycle operations & 16 scenarios on two disposable libraries & create, move, delete and forced re-index pass; two failures (task status, update check) & E2\\
Sub-library hierarchy & Three-level disposable fixture & create, parent-to-child move, listing, reading and judging all pass & E2\\
Tags and graph & 316-tag registry; 658-node graph & 165 stored parts pass the tag check in 47.6\,s; four tag edges on one document move 4{,}675 to 4{,}679 & E2\\
Experience store & One small knowledge base; one end-to-end scenario & 13 candidate drafts; one lesson retrieved first at 0.711; approval path unexercised & E3\\
\bottomrule
\end{tabular}
\end{table*}

\begin{table*}[t]
\centering
\footnotesize
\setlength{\tabcolsep}{4pt}
\caption{Design and headline results of the retrieval and answer experiments.
Each row names the corpus and scale, the metric and the comparison it is read
against, the headline outcome, and the claim that outcome supports together
with its limit. Numbers are those of the cited subsections and their frozen
artefacts; none is re-estimated.}
\label{tab:retrieval-design}
\begin{tabular}{@{}>{\raggedright\arraybackslash}p{0.10\textwidth}>{\raggedright\arraybackslash}p{0.20\textwidth}>{\raggedright\arraybackslash}p{0.24\textwidth}>{\raggedright\arraybackslash}p{0.24\textwidth}>{\raggedright\arraybackslash}p{0.13\textwidth}@{}}
\toprule
Exp & Corpus and scale & Metric and comparison & Headline result & Supports / limit\\
\midrule
E4 & Industrial 42 docs, 16 q; arXiv 79 papers, 10 q & Hit@5, Recall@5, MRR@5 over nine configurations & BM25 and RRF 1.00 vs HMRS-Hybrid 0.44 (industrial); all 0.90 except Graph 0.20 (arXiv) & No consistent hybrid gain; small corpora, retired wrapper\\
E4 verdict & Same logged pools, 191 scores & Fusion order vs judge-as-ranker vs judge-as-filter & 0.94 / 0.44 / 0.94; sign test $p=0.008$ on eight discordant questions & Ordering, not filtering, is the loss; threshold rejects none\\
E2c & Ten-domain 72 q; content-layer arms & Containment; tiered vs flat vs precise & A0h 72.2\% vs A1b 54.2\% (13--0, $p=0.00024$); vs PD 63.9\% (7--1, $p=0.0703$) & Tiered beats flat; navigation-over-precise directional; uncorrected family\\
E5 & Two 30-cell grids, three interfaces & Gold-ID, lexical coverage, citation markers & Gold-ID 60/60; lexical and marker counts vary with substrate & Traceability, not semantic support; one run per arm, at ceiling\\
E5b & Three arms $\times$ three questions, public API & Cost, refusal behaviour, transport & S2 false refusal on a 131-doc catalogue (74.9\,KB over tool budget); repaired by pagination & Transport and cost boundary; one campaign plus replication; per-arm costs in supplement\\
E6 & 191 logged scores; second NLI signal & Gold vs non-gold score separation; rank replay & Gold median 0.88 vs non-gold 0.86, $p=0.26$; NLI Hit@5 0.3125 & The judge is not a ranker; one corpus\\
E7 & Concurrency sweep, 90/180/270 requests & Wall clock, timeout rate at 5/15/30 clients & 41\% timeouts at 30 clients; p50 5.2 to 1577\,ms & Single-instance envelope; not validated throughput\\
E8 & Gate on/off, $n=10$; pinned vs unscoped scope & Answer grades; Hit@5 under each scope & Gate effect null; both scopes Hit@5 1.0 & No answer-quality benefit here; scope not isolated, filename confound\\
\bottomrule
\end{tabular}
\end{table*}

\subsection{E1: Acquisition and parsing}
\label{sec:e1}

The acquisition campaign targeted 100 submitted papers for content-informed
routing into five domain libraries; its campaign record contains 322 stored
parts after the 30{,}000-character split. Storage and retrievability diverge:
at the 2026-10-13 evaluation snapshot, 91 names appear
in the ingest state file and 79 papers were retrievable, so twelve recorded
papers never became retrievable, and nine submitted papers never entered the
state file. The 21 unretrievable papers are the sum of those twelve and nine,
and neither attempts against failed papers nor per-paper parse status is
recorded, so the 21 cannot be attributed to parsing alone
\cite{missier-prov-2013}. Reported failure modes are formula-dense scans and
one tool-layer crash during batch parsing, restarted from a checkpoint.
Table~\ref{tab:kb-checks} lists every functional check with its scale and
outcome.

The cross-lingual benchmark ran twice over 400 pages per language and
reproduced round~1 bit-for-bit, with ingestion success 1.0. Table~\ref{tab:retrieval} reports
storage completeness, an index self-check at Hit@1, and membership figures;
the self-check verifies post-ingestion index reachability, not task-level
retrieval quality. Its English membership accuracy of 0.875 is an organisation
defect---the page stored in the wrong library roughly one time in eight---not
a parsing defect. Membership is scored against the benchmark's source category labels, not a
naming convention derived by the ingest agent, so it is the one
filename-independent placement measurement in this paper.

A second 50-paper, 26-domain campaign parsed and verified every paper, indexed
them in batch, reported all probes ok, and left 41 of 50 documents with a
negative tag flag even though verification passed, so tag-quality failures
never blocked ingestion.
One stage re-ran on inbox ordering, and more than 800 extracted images were
not uploaded. Descriptions and tags are checked against content before
storage~\cite{park-metadata-quality-2009}.


\subsection{E2: Organisation, consistency, and management operations}
\label{sec:e2}


Organizational correctness is three-source consistency: disk files, the tree
index and per-library metadata must agree on document existence, location and
indexed content. Under that definition we found one real defect in a
disposable test library: after a rename the vector index kept the old name's
chunks---count 18 where 9 were expected, a duplicate search matching the
document against its own former path at 0.9738 similarity, a vector probe
returning every chunk twice at identical scores. Only an explicit force
re-index removed them, restoring the count to 9. The failure is invisible from
the metadata alone.

The reorganisation campaign over the 322-part corpus passed its content
check on all 322 stored parts and its three-source consistency check on a
33-part sample, one tenth of the corpus; the duplicate review held no true
duplicate and deleted nothing. Sub-library splitting and hierarchy
creation ran on a disposable three-level fixture: creation with a parent
reference, a parent-to-sub-library move, and listing, reading and judging
through sub-library addresses all returned correct results. A parent catalogue
lists its sub-libraries as container entries carrying their own descriptions,
the one-call view that ladder navigation reads.

Growing from 50 to 100 papers cost no retrieval quality: dense top-1 hits moved
from 48/50 to 49/50 and MRR from 0.98 to 1.00, recorded as a post hoc
held-within-5-pp band
$\pm$5\,pp. In the 16 lifecycle scenarios, create, move, delete and forced
re-index behaved as specified (Table~\ref{tab:functions}). Two failures remain: the task-status
tool reports \texttt{running} after \texttt{finished\_at} is written, and the
update check timed out after 30\,s while every other probe passed. Merge and
split were never exercised under controlled conditions, so we claim nothing
about them.

The tag registry held 316 tags; the tag-quality check passed 165
stored parts with zero failures in 47.6\,s, and orphan tags accumulate (test tags,
section titles), for which a dry-run detector exists. In one case test-library
tags leaked into the main registry, the failure mode vocabulary-based
organisation is prone to~\cite{golder-tagging-2006}. The document graph is
navigation, never evidence: four tag edges on one document moved the edge count
from 4{,}675 to 4{,}679.

\subsection{E2b: Paired route-card audit}
\label{sec:routing-audit}

A separate audit run on 2--3 November 2026 used a ten-domain corpus containing
138 documents in 879 stored parts. The paired question set contains 72
in-domain gold questions; each arm also included four out-of-domain probes,
and two cross-placeholder probes were excluded. Route correctness is the
binary match between a scripted route result and the question's gold-domain
label, measured before document retrieval or answer generation. A0 returned
the labelled domain for 72/72 questions (100\%; Wilson 95\% CI [0.949,
1.000]); the nominal name-only A3 condition did so for 37/72 (51.4\%; CI
[0.401, 0.626]). The paired table has 35 A0-only correct routes and no A3-only
correct routes; exact McNemar $p=5.82\times10^{-11}$ (Figure~\ref{fig:routing}).

This was a scripted route-card probe, not an LLM agent selecting a library.
The route-card vector index is a diagnostic stand-in for the collection tier;
in the deployed system the collection choice is made by the ingest and answer
workflows, so the audit measures the information content of the tier rather
than the production selector. The audit therefore measures the collection tier of the access path in
isolation: it asks which collection a question belongs to, before any document
is retrieved and before any answer is generated. The runner first queried a route-card vector index and, when that returned no
hit, fell back to a deterministic name-and-description matcher over the full
catalogue. In A3, the route cards were created without descriptions, but the
fallback could still read the original descriptions; traces do not record
whether the fallback ran. Thus, the observed A0/A3 difference is not an
isolated causal estimate of description removal. Route correctness and the
separate downstream evidence-pack outcome are reported as distinct metrics.

\begin{figure*}[t]
\centering
\includegraphics[width=\textwidth]{fig3_routing}
\caption{Paired scripted route audit (72 in-domain questions). (a) Gold-domain
choice: A0 72/72, A3 37/72 (Wilson 95\% intervals; McNemar
$p=5.82\times10^{-11}$). (b) Exact-string containment: A0/A3 Laya evidence
packs and A1 concatenated text from all 879 parts (37/72, 33/72, 72/72; Wilson
intervals; paired $p=0.4545$ for A0--A3 and $p=5.82\times10^{-11}$ for
A1--A0). (c) Trace-recorded text bytes on a symmetric-log axis linear through
1 KB, including 11 zero-byte A3 cases; (d) tool wall-clock on a log axis.
Boxes show IQR, whiskers 5th--95th percentiles, dots questions and diamonds
means. A0/A3 bytes sum candidate-snippet and evidence-pack text (possible
overlap); A1 sums document-read responses. Tool time excludes answer
generation. A1 is a script, not an answering agent; the A3 fallback is
unlogged. These are routing, substring and tool-cost diagnostics, not answer
quality or an isolated description effect.}
\label{fig:routing}
\end{figure*}

The configured route-card arms differ on this question set, but the unlogged
A3 fallback prevents attributing that difference solely to descriptions; the
result does not establish general placement or end-to-end answer quality. In
the same run, exact strings occurred in the Laya evidence pack for 37/72 A0
questions and 33/72 A3 questions (paired $p=0.4545$). The A1 read-all script
checked the concatenated text returned from all 879 parts and found 72/72; it
averaged 24.3 MB of trace-recorded text and 253.7 s across 897 tool calls per
question, versus 31.4 KB and 16.3 s across four calls for A0. A3 averaged
23.5 KB and 7.8 s across 4.1 calls. For A0/A3 the byte count sums top-ranked
candidate snippets and evidence-pack text, which may overlap; for A1 it sums
the document-read responses. These counts are not network traffic or measured
LLM context, and tool time excludes answer generation. A1 did not generate
answers, so this is a mechanical full-corpus containment upper bound, not a
careful-agent retrieval or QA baseline. An audit found
exact strings for 14 anchors in their source PDFs, but did not verify those
strings in parser output or indexed parts; containment therefore does not
validate the full ingestion-to-evidence path. The flat and no-gate arms
(A1b, A2) were re-run on 2026-11-10 as scoped per-library directed searches
merged by score, replacing an earlier unscoped protocol; they are reported
with the content-layer optimisation arms in Section~\ref{sec:ladder-opt}.
The question set is unbalanced across
domains (52/72 questions were automatically expanded, with 47/72 drawn from
three domains), so it does not estimate performance on a balanced domain
sample. The expansion seeds and procedure are not recorded in the frozen
artefacts; we therefore report 72/72 as a scripted-probe configuration
comparison, not as a placement-accuracy estimate, and flag the possible
leakage of gold-domain labels into question construction as a limitation
(Section~\ref{sec:threats}). These limitations are discussed in Section~\ref{sec:threats}; raw
traces and aggregate statistics are in
\path{review/experiment-redesign-20261101/artifacts/final_bench/s4_traces_final.json}
and \path{review/experiment-redesign-20261101/artifacts/final_bench/s6_stats.json}.

\subsection{E2c: Content-layer optimisation and the ladder-versus-precise
contrast (2026-11-11)}
\label{sec:ladder-opt}

The route audit isolates collection selection; it does not say whether the
tiered path, once a collection is chosen, retrieves better than a flat
alternative. A follow-up campaign on the same 72 in-domain questions therefore
varied the content layer on the tiered route (A0f, A0h) and added a precise
per-document arm that abandons navigation (PD); the head-to-head
A0h-versus-PD contrast therefore involves both the content and the routing
layer, while the flat and ablated controls (A1b, A2, A3) hold the scope fixed
(Figure~\ref{fig:ladder}). All rates
are gold-string containment over the 72 questions, with Wilson 95\% intervals
and exact paired McNemar tests against the tiered baseline A0
(\path{review/experiment-redesign-20261101/artifacts/final_bench/s6_stats.json}).

The tiered baseline A0 contained the gold string for 37/72 (51.4\%). Letting
the judge read the full document window (A0f) raised this to 45/72 (62.5\%,
$p=0.0078$); adding a vector chunk-window union to the judge pack (A0h)
raised it further to 52/72 (72.2\%, $p=6.1\times10^{-5}$). The precise
per-document arm PD, which searches the ten libraries as a directed union and
judges full windows without navigation, reached 46/72 (63.9\%,
$p=0.0039$). The scoped flat union without a judge (A1b) reached 39/72
(54.2\%, $p=0.6875$), the no-gate arm (A2) 25/72 (34.7\%, $p=0.0227$), and the
name-only card arm (A3) 33/72 (45.8\%, $p=0.4545$); the mechanical read-all
upper bound contained 72/72.

The same-corpus contrasts carry the management claim. Against the scoped flat
union without a judge (A1b), the tiered path with the chunk-window union
(A0h) differs on thirteen questions, all A0h-only (exact McNemar
$p=0.00024$); the precise arm PD also exceeds A1b (7--0, $p=0.0156$). The
head-to-head A0h-versus-PD contrast, which probes the value of navigation on
top of precise per-document judging, differs on eight questions (seven
A0h-only, one PD-only; exact McNemar $p=0.0703$), a consistent direction that
does not reach the 0.05 threshold on this sample. We therefore report the
tiered path as numerically higher than abandoning navigation, not as
equivalent or superior: an equivalence claim would require a margin and power
analysis we did not run. Cost is reported on both dimensions: A0h issued 10.7
search calls and read 309\,KB per question versus 17.0 calls and 308\,KB for
PD (trace means; Figure~\ref{fig:ladder}(b)), so the navigation arm saves
calls but not judge-input bytes. The paired tests are reported uncorrected;
under Holm correction across the nine pairwise contrasts the
tiered-versus-baseline and tiered-versus-flat comparisons
($p=6.1\times10^{-5}$, $0.00024$, $0.0039$, $0.0078$) remain significant, while
the no-gate contrast (A2) and the A0h--PD head-to-head do not. Because A1b is
flat and judge-free, the A0h--A1b gap reflects the combined content-layer
additions (judging plus the window union) rather than tiering alone; the
navigation-isolating contrast is A0h--PD. As an external
reference point, a separate 40-question floor probe over the same libraries
measured gold-document hit rates of 60.0\% (BM25), 47.5\% (dense) and 65.0\%
(RRF) (\path{review/experiment-redesign-20261101/artifacts/final_bench/s5_floor.json});
the ladder arms' containment is thus comparable to simple fusion without
navigation, on a different metric and question set.
An attribution of the 35 A0 misses locates five in window truncation and
thirty in judge or candidate-selection failure; A0f repairs eight, A0h
fifteen and PD nine, so the chunk-window union recovers misses that full-window
judging alone cannot. These are containment diagnostics on one corpus, not
answer-quality or generalisation claims.

\begin{figure*}[t]
\centering
\includegraphics[width=\textwidth]{fig_ladder_optimization}
\caption{Content-layer optimisation arms on the 72 in-domain questions of the
ten-domain corpus (2026-11-11). (a)~Bars are gold-string containment; error bars
are Wilson 95\% intervals; the dashed line is the tiered baseline A0. A0f
(full-window judge) and A0h (plus a chunk-window union) improve on A0; PD
(precise per-document, no navigation) is the ladder-versus-precise contrast
(A0h vs PD: 7--1 discordant, McNemar $p=0.0703$); A0h vs the scoped flat union
A1b is 13--0 ($p=0.00024$). A1b, A2 (no gate) and A3 (name-only cards) are the
flat and ablated controls; read-all is a mechanical upper bound, not an agent
baseline. (b)~Per-question cost on both dimensions: mean search calls and mean
trace-recorded kilobytes (log axis); A0h saves calls relative to PD but reads
a similar byte budget.}
\label{fig:ladder}
\end{figure*}
